% 5.4 章 · 不确定性量化与全局敏感性
% 这一章的笔记要把"给一个不确定的输入、问输出的分布与它的敏感度"拆成四件互不相同的事：
%   1) 蒙特卡洛：误差标度 1/sqrt(N) 到底保证了什么、没保证什么（覆盖率、重尾、误差估计）；
%   2) 拟蒙特卡洛：低偏差点集把误差压下去多少，代价是什么（网性质、置乱、没有误差棒）；
%   3) 全局敏感性：Sobol' 方差分解在什么前提下成立，前提没了会发生什么（相关输入不唯一）；
%   4) 多项式混沌：把模型换成一个谱展开的代理，什么时候能换、什么时候绝不能外推。
% 所有结论都配一条带实测数字的断言，附表由 tools/build_book.py 生成后 \input 进来。

\numlabchapter{不确定性量化与全局敏感性}{5.4}

\section{这一章要解决什么}

输入不是数而是分布时，"算一次"就不再是问题的全部。本章把下面五句话算清楚：

\begin{enumerate}
  \item \textbf{蒙特卡洛的误差是 1/sqrt(N)，这句话保证什么？} 它保证\emph{标度}：
        $N=64\to4096$ 上实测双对数斜率 $0.4731$（$1/N$ 轴，理论 $0.5$）。
        它不保证覆盖率：$N=100$ 时正态输出实测 $0.9533$，重尾输出只有 $0.9267$
        （差 $+0.0267$，二项标准误 $0.0122$）；$N=1000$ 时两者回到 $0.9467/0.9567$
        （定理~\ref{thm:mc}）。
  \item \textbf{"误差棒"是从哪来的？} 只有随机化才有误差棒。确定性 Sobol' 点集在
        $N=4096$ 上的复制标准差实测 $0.0$，而真实误差是 $2.866\times10^{-4}$——
        覆盖率 $0.00$；同一个点集线性置乱后标准误 $6.997\times10^{-7}$，
        $R=64$ 次复制的"诚实度" $\text{bar}/\text{rmse}\cdot\sqrt R=0.986$
        （定理~\ref{thm:qmc}）。
  \item \textbf{低偏差点集能快多少？} 帧所用窗口 $N\in[2^6,2^{12}]$ 上，
        MC 斜率 $-0.4969$、随机化 Sobol' $-1.2185$（差 $0.7216$）；
        换到 $[2^{10},2^{16}]$ 差距拉到 $1.3883$——\emph{斜率本身依赖窗口}。
        而 $1[x_1<1/2\wedge x_2<1/2]$ 这个间断函数被 Sobol' 网\emph{精确}积分：
        $14$ 组 $(N,\text{随机化})$ 的误差全是 $0$；把边界挪到 $0.3$ 就回到 $-0.898$
        （定理~\ref{thm:qmc}）。
  \item \textbf{Sobol' 指数是恒等式还是判据？} 独立输入下它是判据：
        g-函数实测 $\sum_iS_i=0.9348\pm0.0820$（解析 $0.9290$）、
        $\sum_iS_{Ti}=1.1077\pm0.0358$（解析 $1.0716$），$1$ 被夹在中间。
        输入相关时前提没了：双线性模型 $\rho=0.5$ 时条件方差口径给出
        $\sum_i\mathrm{Var}(E[Y\mid X_i])/V=1.4942>1$，
        同一份样本打散成乘积测度口径给 $0.9890$（定理~\ref{thm:sobol}）。
  \item \textbf{代理模型什么时候能信？} 多项式混沌在光滑函数上是谱收敛：
        $e^{1.5x}$ 的 $\log_{10}\text{err}$ 对阶数 $p$ 的斜率实测 $-1.0005$
        （$R^2=0.992$），$p=1\to12$ 掉 $10.9$ 个数量级；
        折点函数 $\lvert x\rvert$ 只有代数衰减（实测指数 $1.094\sim1.306$）。
        $n_{\text{train}}=8$、$p=10$ 时训练误差 $4.49\times10^{-16}$ 而留出误差
        $4.855\times10^{-1}$——差 $10^{15}$ 倍；$2$ 倍外推放大 $709$ 倍（$e^{1.5x}$）
        与 $1.4985\times10^4$ 倍（Runge）（定理~\ref{thm:pce}）。
\end{enumerate}

记号：$X\in\mathbb{R}^d$ 是输入随机向量，$Y=f(X)$ 是输出，$V=\mathrm{Var}(Y)$；
$X_A=(X_i)_{i\in A}$；$\mu_A=\mathbb{E}[Y\mid X_A]$、$V_A=\mathrm{Var}(\mu_A)$；
$S_A=V_A/V$、$S_{Ti}=1-V_{-i}/V$；
$\{\phi_\alpha\}$ 是多重指标 $\alpha$ 下的 gPC 基；
$D_N^\ast$ 是星偏差，$V_{\mathrm{HK}}(f)$ 是 Hardy--Krause 意义下的变差；
$N$ 是样本数、$R$ 是随机化复制数、$p$ 是截断阶数。

生成元：本章所有随机数只经 \texttt{numcore.utils.rng}（P4），
$\rho=0.5$ 的相关样本用 Cholesky 因子 $L$ 从独立正态造出（$x=Lz$，实测相关系数
$\hat\rho=0.4990$）；Sobol' 点集由方向向量 $\{v_k\}$ 与
\texttt{numcore.linalg.eigen} 的 Jacobi 特征分解（Golub--Welsch 造 Gauss 规则）驱动，
不使用任何现成的低差异序列或正交多项式库。

\section{蒙特卡洛：标度与保证是两件事}\label{sec:mc}

\begin{theorem}[蒙特卡洛误差的标度与它的三个前提]\label{thm:mc}
设 $Y=f(X)$、$\sigma^2=\mathrm{Var}(Y)<\infty$，$\hat\mu_N=\frac1N\sum_{k=1}^NY_k$ 是
$N$ 个独立样本上的均值。则 $\mathbb{E}[\hat\mu_N]=\mu$、
$\mathrm{RMSE}(\hat\mu_N)=\sigma/\sqrt N$；且中心极限定理给出
$\sqrt N(\hat\mu_N-\mu)/\sigma\Rightarrow\mathcal N(0,1)$。

三条推论各有边界：
\begin{enumerate}
  \item \textbf{标度}：$\log\text{RMSE}$ 对 $\log N$ 的斜率是 $-1/2$。
        实测（$N=2^6,\dots,2^{13}$，$R=32$）$0.4731$，与
        \texttt{numcore.verify.fit\_order} 的口径一致。
  \item \textbf{覆盖率}：$\hat\mu_N\pm z_{0.975}\hat\sigma/\sqrt N$ 的覆盖率只在
        "$\hat\sigma$ 稳定 + 分布不太歪"时接近 $0.95$。$N=100$ 时重尾输出实测
        $0.9267$（半宽 $0.2856$），正态输出 $0.9533$（半宽 $0.2781$）。
  \item \textbf{有限方差}：$\sigma=\infty$ 时上述全部失效；这也是后面
        "间断/重尾函数上 QMC 的优势会消失"的同一根源。
\end{enumerate}
\end{theorem}

\begin{proof}
前两条是教科书结论：$\mathrm{Var}(\hat\mu_N)=\sigma^2/N$ 来自独立性，
CLT 由 Lindeberg--Levy 定理。第三条：若 $\mathbb{E}\lvert Y\rvert^2=\infty$，
则 $\hat\sigma^2$ 不收敛，覆盖率的定义本身就不稳定。实测部分见
\texttt{tests/5.4/test\_mc\_confidence.py}：300 次独立实验里，$N=100$ 的正态输出
命中 $286/300$、重尾输出命中 $278/300$，差 $8$ 次 $=2.67\%$，是二项标准误
$1.22\%$ 的 $2.2$ 倍——"$1/\sqrt N$ 是标度不是保证"这句话在这里是可测量的，
而不是修辞。
\end{proof}

\paragraph{四个可以直接看的量}
（1）\textbf{分布本身}：均匀/正态/Beta/指数四种采样器的一阶二阶矩实测偏差
$\le0.0023$。（2）\textbf{输出的密度}：均匀输入 + 线性模型有\emph{精确}的三角密度
（卷积），经验密度与它的 $L_1$ 距离实测 $0.01402$、质量 $0.999981$；
指数输入没有解析卷积，函数返回 \texttt{None} 而不是编一个"参考密度"。
（3）\textbf{直方图}：密度积分 $1.000000000000$、计数恰为 $N$，
bin 数的误差曲线在 $B=16$ 取极小（$0.0029$），两端 $0.0052$（$B=8$）与
$0.0062$（$B=96$）——偏差-方差权衡的形状。
（4）\textbf{KS 距离}：对\emph{正确}的参考线，KS 随 $N$ 下降
（$0.01492\to0.00722\to0.00422$）；对 CLT 近似参考线，它停在
$0.01766$ 的地板上不再下降——KS 小只说明"和这条参考线接近"，
不说明参考线是对的。

\section{拟蒙特卡洛：低偏差、网性质与"没有误差棒"}\label{sec:qmc}

\begin{theorem}[Koksma--Hlawka 与 $(t,m,s)$-网的精确性]\label{thm:qmc}
设 $f$ 在 $[0,1)^d$ 上有界变差（Hardy--Krause），$P_N$ 是 $N$ 点集，则
\[
  \Bigl\lvert \frac1N\sum_{k=1}^N f(x_k)-\int_{[0,1)^d}f\Bigr\rvert
  \;\le\; V_{\mathrm{HK}}(f)\,D_N^\ast(P_N).
\]
若 $P_N$ 是 $2$ 维 $(0,m,2)$-网（$N=2^m$），则对每个
$m_1+m_2=m$，任一初等区间
$[a_12^{-m_1},(a_1+1)2^{-m_1})\times[a_22^{-m_2},(a_2+1)2^{-m_2})$
里\emph{恰好}有一个点；因此任何"轴对齐、边界在 $2^{-m}$ 网格上"的矩形指示函数
被\emph{精确}积分。
\end{theorem}

\begin{proof}
不等式是 Koksma--Hlawka 不等式（$V_{\mathrm{HK}}$ 为变差、$D_N^\ast$ 为星偏差）。
网的精确性：把点 $x_k$ 的前 $m_1$ 位与 $x_k$ 的前 $m_2$ 位拼成一个 $m$ 位字，
网性质说这个映射是双射；于是每个初等区间恰好命中一个点，
矩形指示函数的均值 $=$ 它的体积。
\end{proof}

\paragraph{实测：不平等式与"上界不是估计"}
$5$ 个被积函数 $\times$ $4$ 种点集 $=20$ 组，不等式一次都没被违反；
但松紧度 $\text{err}/\bigl(V_{\mathrm{HK}}D_N^\ast\bigr)$ 从 $0.000$ 跨到 $0.968$
（$3$ 个数量级）——同一个上界，有的函数上比真实误差大 $10^3$ 倍以上。
变差本身也要小心：$\sin(2\pi x)\sin(2\pi y)$ 的 $k$ 倍频版本实测满足
$V(k)/(V(1)k^2)=1$ 到 $9.3\times10^{-4}$（$400$ 格差分），
网格加密到 $1600$ 后偏差再降一半以上——偏差来自差分，不是标度律错了。

\paragraph{实测：斜率、窗口与"对齐的陷阱"}
\begin{center}
\begin{tabular}{lrrr}
\hline
窗口 & MC & 随机化 Sobol' & 随机化格点 \\
\hline
$[2^{6},2^{12}]$（帧所用） & $-0.4969$ & $-1.2185$ & $-0.730$ \\
$[2^{10},2^{16}]$ & $-0.5151$ & $-1.9034$ & $-0.264$ \\
\hline
\end{tabular}
\end{center}
同一批被积函数与随机化方式，只换窗口，Sobol' 的实测斜率从 $-1.22$ 变到 $-1.90$。
这不是测量误差，而是对数因子与网对齐效应在不同 $N$ 区间权重不同。
\textbf{任何"实测收敛率"都必须连窗口一起引用}，否则就是拿窗口当结论。
格点在 $[2^{10},2^{16}]$ 上退化（$-0.264$）也是同一件事的另一面：
它只在生成向量设计的那段 $N$ 上才表现出格点的优势。

间断函数的对照更能说明"网性质"与"快"不是一回事：

\begin{center}
\begin{tabular}{lrr}
\hline
被积函数 & MC 斜率 & Sobol' 斜率 \\
\hline
$1[x_1<1/2\wedge x_2<1/2]$（边界在 $2^{-1}$） & $-0.541$ & 测不出（误差恒为 $0$）\\
$1[x_1<3/10\wedge x_2<3/10]$（边界不在 $2^{-m}$） & $-0.542$ & $-0.898$ \\
\hline
\end{tabular}
\end{center}

第一行里 $14$ 组 $(N,\text{随机化})$ 的误差全是 $0.0$——因为 $0.5$ 是 $2$ 的负幂，
$[0,1/2)^2$ 是两个初等区间的并，网性质直接给出精确体积。
帧必须把这种情形报成"误差精确为 $0$、斜率不可测"（\texttt{error\_valid} 全 $0$、
\texttt{slope\_current=None}），否则读者会把结构性精确读成"QMC 快得离谱"。

\paragraph{没有随机化就没有误差棒}
确定性 Sobol' 点集在 $N=4096$ 上的 $R=64$ 次"复制"逐位相同：标准差 $0.0$，
而真实误差 $2.866\times10^{-4}$，于是"$95\%$ 置信区间覆盖率"实测 $0.00$。
线性置乱（Matousek）之后标准误 $6.997\times10^{-7}$，与真实 RMSE 的比
$\text{bar}\cdot\sqrt R/\text{rmse}=0.986$——误差棒这才真的是误差棒。
批量均值法是另一条路：它与真实标准误的比实测 MC $1.143$、Sobol' $0.810$、
格点 $0.954$；QMC 那一侧系统性偏小（分批之间负相关），所以本章报误差棒时
优先用"独立随机化复制"。

\paragraph{线性置乱必须同时满足两件事}
方向向量的数字置乱若写成"输出第 $b'$ 位由输入更细的位组合"或者"对角线不钉成 $1$"，
一维分层看起来仍然完美（每格计数 $[16,16]$），但二维网性质已经没了：
实测 $m=4$ 时格子计数从"恰好 $1$"变成 $[0,4]$，
用 GF(2) 秩判据可以逐格验出失败。修好之后，$m=4/8/10$、两种种子、
三种置乱（无/数字移位/线性）在每个 $m_1+m_2=m$ 的切法上都恰好一个点。
这条写进断言是因为它**只能**这样测出来：只看边缘分布永远发现不了。

\section{全局敏感性：判据不是恒等式}\label{sec:sobol}

\begin{theorem}[独立输入下的 Sobol' 分解]\label{thm:sobol}
设 $X_1,\dots,X_d$ 相互独立、$f\in L^2$，则 $f$ 有唯一的正交分解
$f=f_0+\sum_i f_i(X_i)+\sum_{i<j}f_{ij}(X_i,X_j)+\cdots$，
且
\[
  V=\sum_{A\ne\varnothing}V_A,\qquad
  \sum_i S_i\le1\le\sum_i S_{Ti},\qquad
  S_{Ti}=S_i+\sum_{A\ni i,\,A\ne\{i\}}S_A .
\]
两个不等式都是判据：两侧的差恰好是交互项之和
$\sum_{\lvert A\rvert\ge2}S_A$。
\end{theorem}

\begin{proof}
独立性给出 $L^2(\prod\mu_i)$ 的张量积正交分解，逐项取条件期望即得唯一性；
对 $V=\mathrm{Var}(f)$ 用正交性展开得 $V=\sum_A V_A$。
把 $S_i$ 求和：每个 $\lvert A\rvert=1$ 的项算一次、$\lvert A\rvert\ge2$ 的项一次不算，
故 $\sum_iS_i=V^{-1}\sum_{\lvert A\rvert=1}V_A\le1$；
而 $S_{Ti}=V^{-1}\sum_{A\ni i}V_A$，每个 $\lvert A\rvert\ge1$ 的项被算
$\lvert A\rvert\ge1$ 次，故 $\sum_iS_{Ti}\ge1$；
两式相减即得 $V^{-1}\sum_{\lvert A\rvert\ge2}V_A$。
\end{proof}

\paragraph{实测：两个经典函数}
Ishigami（$a=7,b=0.1$）闭式：$S=(0.3139,0.4424,0.0000)$、
$S_T=(0.5576,0.4424,0.2437)$，$\sum S_i=0.7563$、$\sum S_T=1.2437$——
交互项 $S_{13}=0.2437$ 比第一维的一阶效应还大。
g-函数（$a=(0,1,4.5)$）：$\sum S_i=0.9290$、$\sum S_T=1.0716$；
$N=4096$ 的 pick-freeze（Saltelli/Jansen 两套估计器都在同一帧里）实测
$\sum S_i=0.9348\pm0.0820$、$\sum S_T=1.1077\pm0.0358$，$1$ 被区间夹住。
顺便说：**一阶指数之和大于 $1$ 是"代码写错了"的判据**——
本模块第一版把 $V_j$ 写成 $\sigma_j^2\prod_{k\ne j}(1+\sigma_k^2)$，
给出 $\sum_jS_j=1.0716>1$，看起来完全正常；正确写法是 $0.9290$。

\paragraph{实测：估计器的漂移是偏差，不是噪声}
同一函数、同一随机流，只取前 $N$ 行：$\sum_iS_i(N)$ 实测
$1.1587,0.8416,1.1657,1.0034,0.8684,0.9320,0.9488$
（$N=64,128,256,512,1024,2048,4096$），漂移幅度 $0.3242$，
末端与解析值差 $0.0198$。小 $N$ 处的 bootstrap 半宽能到 $\pm0.35$，
所以单个 $N$ 上"反着走"完全是可能的——**这正是它必须被画出来的原因**：
把 $N$ 从小拖到大，读者看到的不是"误差变小"，而是"偏差在消失"。

\paragraph{相关输入：分解不再唯一}
$\rho=0.5$ 的双线性模型 $Y=X_1+X_2$（$X\sim\mathcal N(0,\Sigma)$）：
按条件方差口径 $\sum_i\mathrm{Var}(E[Y\mid X_i])/V=1.4942$（超出 $0.4942$），
把同一份样本按列打散成乘积测度后是 $0.9890$。两个数都对，
但它们回答的不是同一个问题：前者问"已知 $X_i$ 能解释多少方差"（在相关设计下
重复计入共享的部分），后者问"若把输入强行解耦，独立贡献是多少"。
本章的帧必须同时给出这两个数 + \texttt{corr\_valid} 掩码 + 文字说明，
否则"$\sum S_i\le1$"会被当成恒等式用。

\section{多项式混沌：谱收敛、代理与绝不外推}\label{sec:pce}

\begin{theorem}[gPC 投影、矩与方差分解]\label{thm:pce}
设输入测度有正交归一基 $\{\phi_\alpha\}$（$\mathbb{E}[\phi_\alpha\phi_\beta]=\delta_{\alpha\beta}$），
$f=\sum_\alpha c_\alpha\phi_\alpha$、$c_\alpha=\mathbb{E}[f\phi_\alpha]$。则
\[
  \mathbb{E}[f]=c_0,\qquad
  \mathrm{Var}(f)=\sum_{\alpha\ne0}c_\alpha^2,\qquad
  S_i=\frac{1}{V}\sum_{\alpha\in\mathcal A_i}c_\alpha^2,\qquad
  \mathcal A_i=\{\alpha:\alpha_i>0,\ \alpha_j=0\ (j\ne i)\}.
\]
按 $L^2$ 截断到 $p$ 阶时，误差 $\lVert f-f_p\rVert_2^2=\sum_{\lvert\alpha\rvert>p}c_\alpha^2$。
\end{theorem}

\begin{proof}
第一式是 $\alpha=0$ 的系数；第二式用 Parseval；
第三式把 $\alpha\ne0$ 按"哪些分量非零"分组，$\mathcal A_i$ 恰好是只依赖 $X_i$ 的项，
由定理~\ref{thm:sobol} 的独立性前提即得。截断误差是同一 Parseval 等式的尾部。
\end{proof}

\paragraph{实测：正交性、求积与"矩恒等式"}
Legendre 基在 $\mathrm dx/2$ 下、Hermite 基在标准正态下的 Gram 矩阵偏离单位阵
最大 $8\times10^{-14}$ 与 $4.9\times10^{-11}$；
Gauss--Legendre 节点与 2.4 章的参考实现相对差 $\le5\times10^{-15}$，
权重**恰好**是它的两倍（那一支对 $\mathrm dx$ 归一）；
Gauss--Hermite 三点权重是 $(1/6,2/3,1/6)$。
\begin{quote}
本模块第一版在这里差了一个因子 $2$（给出 $(1/12,\dots,1/12)$），
后果是"正交归一"和"$V=\sum c_\alpha^2$"两条会**同时**错而成体系。
所以这一类断言刻意写成"与 2.4 章对拍 + 三点闭式"两条独立的锁。
\end{quote}
投影与回归在超定情形下必须给同一个答案：四次多项式在 $p\ge4$ 时误差
$\le10^{-12}$（\texttt{exact\_at\_order=4}），在 $p=2,3$ 时是
$10^{-0.36},10^{-1.05}$（明显不为零，否则这条断言什么也没测）。

\paragraph{实测：谱衰减的两种形态}

\begin{center}
\begin{tabular}{lrrr}
\hline
被积函数 & 半对数斜率 & $R^2$ & $p=1\to12$ 降幅 \\
\hline
$e^{1.5x}$ & $-1.0005$ & $0.992$ & $10.9$ 个数量级 \\
$\lvert x\rvert$ & $-0.1080$ & $0.860$ & $1.5$ 个数量级 \\
\hline
\end{tabular}
\end{center}

$\lvert x\rvert$ 的代数指数用三种子样本拟合分别是 $1.2439$（全部阶）、
$1.3061$（奇数阶）、$1.0940$（偶数阶）——偶数阶几乎不降，
误差曲线是"台阶"。渐近理论值 $\sim1$（$c_k\sim k^{-3/2}$），但 $p\le16$
上测不出干净的 $-1$，所以本章只说"与 $p^{-1}$ 的代数衰减一致"。

\paragraph{实测：过拟合区与外推失效}
$n_{\text{train}}=8$、$p=10$：训练误差 $4.49\times10^{-16}$、
留出误差 $4.855\times10^{-1}$（比 $3.69\times10^{4}$ 倍），
\texttt{under\_flag} 从 $n_{\text{basis}}>n_{\text{train}}$ 那一刻起为 $1$；
默认帧（$n_{\text{train}}=64$）里 $e^{1.5x}$ 的 2 倍外推放大 $706$ 倍，
Runge 函数放大 $1.4985\times10^{4}$ 倍。结论不是"外推要小心"，
而是\textbf{多项式混沌对外推没有收敛保证}：多项式在支撑集外增长得比任何
有理函数都快，误差随倍率单调恶化（实测 $1.1/1.5/2.0$ 倍单调）。

\section{四张图分别在看什么}

\paragraph{图 1（\texttt{point\_fill}）点集填充对决}
左图把同一被积函数下的 MC / Sobol' / 格点 / 张量格点点集画出来，
右图是双对数误差曲线 + $1/\sqrt N$ 参考带。
诚实机制有三处：（1）误差棒只有 \texttt{randomize=True} 时存在
（\texttt{band\_valid} 全 $0$ + \texttt{band\_available=False} + 文字说明）；
（2）误差曲线贴地板时斜率报 \texttt{None}；（3）对数轴上的 $0$ 写成地板 $10^{-17}$
并配 \texttt{error\_valid} 掩码。默认帧的斜率为
MC $-0.497$、Sobol' $-1.219$、格点 $-0.730$、张量格点 $-1.005$。

\paragraph{图 2（\texttt{output\_hist}）输出分布随输入分布实时变化}
$Y=X_1^2+0.5X_2$ 上换四种输入分布，直方图 + 精确/近似参考密度 + 副轴 bin 数误差曲线。
参考线的身份必须声明：均匀+线性给的是\emph{精确}三角密度
（KS $0.0029$ vs 正态参考 $0.0190$），指数+线性只给 CLT 近似
（\texttt{ref\_is\_exact=False}，文字写明）——否则 $\lvert$密度差$\rvert$
会被读成"模型误差"。实测默认帧（均匀+二次）偏度 $+0.433$、
超额峰度 $-0.559$、KS（经验 vs CLT）$0.0584$。

\paragraph{图 3（\texttt{sobol\_bars}）Sobol' 指数柱状图}
一阶 vs 总效应柱 + 拖 $\rho$ 看交互浮现。$\rho$ 只对正态域模型（双线性、相关二次）
有定义；均匀域模型（Ishigami、g-函数）上帧报
\texttt{corr\_applicable=False} + 文字说明，而不是编一个相关分解。
$\rho=0.5$ 时同一帧里同时给出三个数：
\texttt{first}（独立 pick-freeze）、\texttt{sum\_between}（条件方差口径 $1.4942$）、
\texttt{sum\_indep}（乘积测度口径 $0.9890$），外加
\texttt{decomposition\_note} 说明它们各自回答什么问题。

\paragraph{图 4（\texttt{pce\_spectrum}）谱系数衰减与代理外推}
左图是 $\log_{10}\lvert c_\alpha\rvert$ 的衰减（含谱衰减率），
右图是阶数-误差曲线（训练误差与留出误差分开画）+ 外推曲线。
过拟合区、\texttt{exact\_at\_order}、\texttt{extrap\_ratio} 都在同一帧里，
让"训练误差到机器零"与"留出误差 $0.485$"同时可见。

\section{本章的"不要被骗"清单}

\begin{itemize}
  \item $1/\sqrt N$ 是\emph{标度}，不是保证：重尾 + 小 $N$ 时覆盖率掉到 $0.9267$。
  \item 误差棒不是"算出来"的：确定性点集给不出误差估计（复制标准差 $0$、
        真实误差 $2.866\times10^{-4}$、覆盖率 $0.00$）。
  \item 实测斜率必须连窗口一起报：同一函数在 $[2^6,2^{12}]$ 与 $[2^{10},2^{16}]$
        上是 $-1.22$ 与 $-1.90$。
  \item 网精确积分对齐的间断函数（误差恒为 $0$）不是"快得离谱"，
        换个边界位置（$0.3$）就回到 $-0.898$。
  \item $\sum_iS_i\le1\le\sum_iS_{Ti}$ 是\emph{独立输入}下的判据：
        $\rho=0.5$ 时条件方差口径给 $1.4942$。
  \item 估计器的漂移（$0.3242$）是偏差：换随机种子不会让它变小。
  \item 训练误差到机器零什么也不能说明：$n_{\text{train}}=8,p=10$ 时留出误差 $0.485$。
  \item 多项式混沌外推没有收敛保证：2 倍外推放大 $706$（$e^{1.5x}$）到
        $1.5\times10^4$（Runge）倍。
\end{itemize}

% 本文件由 tools/build_book.py 从 dist/reports/assertions.json 生成，请勿手改。
\begin{longtable}{@{}p{0.06\textwidth}p{0.32\textwidth}p{0.22\textwidth}p{0.16\textwidth}@{}}
\caption{章节 5.4 数值断言（生成时间 2026-10-18T18:30:41-0400）}\\
\toprule
类别 & 断言 & 理论值 & 实测值 \\
\midrule
\endfirsthead
\toprule
类别 & 断言 & 理论值 & 实测值 \\
\midrule
\endhead
\bottomrule
\endlastfoot

A & 端点组合帧数与 JSON 里真正的非有限常量个数 & 非有限常量个数 =\allowbreak{} 0(\allowbreak{}字符串里出现 NaN 三个字母不算)，帧数 >=\allowbreak{} 60 & 66 帧，计数 0(\allowbreak{}point\_fill: 17 个数组、\allowbreak{}17 个 meta 字段) \checkmark \\
A & 对数曲线范围、地板说明、以及贴地板处掩码确实为 0 & 全部曲线 >=\allowbreak{} -17 且 max > -16；discontinuous 掩码全 0、\allowbreak{}offset\_box 掩码全 1 & point\_fill.error$\in$[-5.5,\allowbreak{}-3.5]；point\_fill.error\_mc$\in$[-2.8,\allowbreak{}-1.8]；point\_fill.error\_sobol$\in$[-5.5,\allowbreak{}-3.5]；point\_fill.error\_ref$\in$[-2.7,\allowbreak{}-1.8]；point\_fill.band\_lo$\in$[-5.6,\allowbreak{}-3.6]；point\_fill.band\_hi$\in$[-5.4,\allowbreak{}-3.4]；output\_hist.error$\in$[-2.9,\allowbreak{}-1.0]；output\_hist.bins\_error$\in$[-1.3,\allowbreak{}-1.3]；pce\_spectrum.error$\in$[-11.5,\allowbreak{}-0.6]；pce\_spectrum.error\_train$\in$[-11.6,\allowbreak{}-0.7]；pce\_spectrum.coeff\_rms$\in$[-3.5,\allowbreak{}0.2]；discontinuous 掩码 [0.]，offset\_box 掩码 [1.] \checkmark \\
A & 三种情形下的 band\_\allowbreak{}valid / band\_\allowbreak{}available / 文字说明 & 关随机化与 grid：全 0 且 band\_available=\allowbreak{}False；开随机化：全 1 且 True & 关：valid=\allowbreak{}[0.]，available=\allowbreak{}False；grid：valid=\allowbreak{}[0.]，available=\allowbreak{}False；开：valid=\allowbreak{}[1.]，available=\allowbreak{}True；说明含「不存在」=\allowbreak{}True \checkmark \\
A & 贴地板时的 slope\_\allowbreak{}current 与说明文字 & slope\_current =\allowbreak{} None；note 里出现「精确为 0」；偏移边界上斜率 $\in$ [-1.2,\allowbreak{}-0.6] & slope\_current=\allowbreak{}None，note 片段=\allowbreak{}'sobol 的误差精确为 0(\allowbreak{}点集与间断/\allowbreak{}边界对齐，见 error\_valid)'，offset\_box 的 slope\_current=\allowbreak{}-0.8698731004959773 \checkmark \\
A & rho=0 与 rho=0.5 的字段对照(双线性模型) & rho=\allowbreak{}0：三个字段为 None、\allowbreak{}corr\_valid 全 0；rho=\allowbreak{}0.5：sum\_between > 1.4 且 sum\_indep $\in$ [0.95,\allowbreak{}1.05] & rho=\allowbreak{}0：corr\_valid=\allowbreak{}[0.]，overlap=\allowbreak{}None，sum\_between=\allowbreak{}None；rho=\allowbreak{}0.5：corr\_valid=\allowbreak{}[1.]，sum\_between=\allowbreak{}1.4942，sum\_indep=\allowbreak{}0.9890，rho\_measured=\allowbreak{}0.4990，corr\_note 带数字=\allowbreak{}True \checkmark \\
A & 均匀输入方体上的模型拖 rho & corr\_applicable=\allowbreak{}False、\allowbreak{}corr\_active=\allowbreak{}False，且文字说明 rho 无效 & corr\_applicable=\allowbreak{}False，corr\_active=\allowbreak{}False，说明含「无效」=\allowbreak{}True \checkmark \\
A & 独立/相关两种情形下的和与判据布尔 & 独立：sum S\_i <=\allowbreak{} 1 <=\allowbreak{} sum S\_Ti 且布尔为真；相关：sum\_between > 1(\allowbreak{}判据在该口径下失效) & 独立 g-函数：sum S\_i=\allowbreak{}0.9348，sum S\_Ti=\allowbreak{}1.1077，lower=\allowbreak{}True，upper=\allowbreak{}True，解析=\allowbreak{}0.9290；相关双线性：sum\_between=\allowbreak{}1.4942(\allowbreak{}>1) \checkmark \\
A & 过拟合区的三个字段 + 训练/留出误差 & underdetermined=\allowbreak{}True；under\_flag =\allowbreak{} (\allowbreak{}n\_basis > 8) 逐点成立；训练 <=\allowbreak{} 1e-14 且留出 > 0.1 & underdetermined=\allowbreak{}True，欠定阶数 5，训练 4.49e-16，留出 4.855e-01，外推 1.392e+03(\allowbreak{}放大 2.87e+03 倍) \checkmark \\
A & poly4 的 exact\_\allowbreak{}at\_\allowbreak{}order / 掩码 / 当前阶误差 & exact\_at\_order =\allowbreak{} 4；p<4 时 log10(\allowbreak{}err) >=\allowbreak{} -1.1；当前阶误差 <=\allowbreak{} 1e-12；掩码全 1 & exact\_at\_order=\allowbreak{}4，log10 误差曲线 [ -0.36  -0.89  -1.05 -14.6  -14.75 -14.38 -14.4  -14.45 -13.73 -13.81
 -13.41 -13.78]，掩码取值 [1.]，当前留出误差 1.651e-14 \checkmark \\
A & 三种输入分布下 ref\_\allowbreak{}is\_\allowbreak{}exact / KS 对照 & 正态+线性：exact=\allowbreak{}True 且两个 KS 相差 <=\allowbreak{} 2e-3(\allowbreak{}精确参考就是正态)；均匀+线性：精确三角密度，ks\_to\_normal >=\allowbreak{} 3 * ks\_to\_ref；指数+线性：没有精确参考，exact=\allowbreak{}False 且 ref\_kind 里写明是 CLT 近似 & 正态：exact=\allowbreak{}True，ks\_ref=\allowbreak{}0.0064，ks\_normal=\allowbreak{}0.0054；均匀：exact=\allowbreak{}True，ks\_ref=\allowbreak{}0.0029，ks\_normal=\allowbreak{}0.0190；指数：exact=\allowbreak{}False，ref\_kind=\allowbreak{}'CLT 正态近似(\allowbreak{}同均值同方差的正态密度，不是真' \checkmark \\
A & 四处数组与核心调用的逐位/对数一致性 & 点集与指数逐位相等(\allowbreak{}差 =\allowbreak{} 0)；计数和 =\allowbreak{} N；误差对数差 <=\allowbreak{} 1e-12 & 点集差 0.0e+00，指数差 0.0e+00，计数和 20000，误差 log10 差 0.000e+00(\allowbreak{}帧的曲线来自同一个 pce.error\_vs\_order 调用) \checkmark \\
A & norm\_\allowbreak{}ppf 与 verify.stats 的最大偏差 + cdf 往返 + 端点行为 & 偏差 =\allowbreak{} 0(\allowbreak{}同一有理逼近)、\allowbreak{}往返 <=\allowbreak{} 1e-8、\allowbreak{}端点为 $\pm$inf & 偏差 0.000e+00，往返 2.465e-10，ppf(\allowbreak{}0)=\allowbreak{}-inf，ppf(\allowbreak{}1)=\allowbreak{}inf \checkmark \\
A & 四种分布的一阶/二阶矩与支撑集 & |mean| <=\allowbreak{} 0.02、\allowbreak{}|std-1| <=\allowbreak{} 0.02 & uniform: mean=\allowbreak{}-0.0009 std=\allowbreak{}1.0006 range=\allowbreak{}[-1.73,\allowbreak{}1.73]；normal: mean=\allowbreak{}+0.0000 std=\allowbreak{}1.0011 range=\allowbreak{}[-4.49,\allowbreak{}4.43]；beta: mean=\allowbreak{}-0.0023 std=\allowbreak{}0.9992 range=\allowbreak{}[-1.79,\allowbreak{}4.08]；exponential: mean=\allowbreak{}-0.0010 std=\allowbreak{}0.9979 range=\allowbreak{}[-1.00,\allowbreak{}12.74]；beta 偏度 0.593(\allowbreak{}Beta(\allowbreak{}2,\allowbreak{}5) 解析偏度 0.596) \checkmark \\
A & 均匀和的经验 L1 距离 / 精确密度积分 / 指数和的返回值 & L1 < 0.02、\allowbreak{}积分 =\allowbreak{} 1(\allowbreak{}误差 <=\allowbreak{} 1e-4 是 60 格黎曼和的求积误差)、\allowbreak{}指数返回 None & L1=\allowbreak{}0.01402，积分=\allowbreak{}0.999981，指数分布的返回值=\allowbreak{}None \checkmark \\
B & 直方图密度积分 / 计入样本数 / bin 误差曲线形状 & 积分 =\allowbreak{} 1(\allowbreak{}1e-12)、\allowbreak{}计数和 =\allowbreak{} N、\allowbreak{}两端误差 > 内部极小 & 积分=\allowbreak{}1.000000000000，计数 50000/\allowbreak{}50000，极小在 B=\allowbreak{}16(\allowbreak{}误差 0.0029)，B=\allowbreak{}8 $\to$ 0.0052，B=\allowbreak{}96 $\to$ 0.0062 \checkmark \\
C & MC 的双对数斜率(框架 fit\_\allowbreak{}order 与本模块 loglog\_\allowbreak{}slope) & 0.5 $\pm$ 0.08，且两条路径一致 & fit\_order 0.4731(\allowbreak{}R$^2$=\allowbreak{}0.987073，n\_used=\allowbreak{}6)，loglog\_slope 0.4731，误差序列 [0.01365 0.01112 0.00848 0.00572 0.00416 0.00268 0.00194 0.00134] \checkmark \\
C & N=100 时两种输出的 95\% 覆盖率与其差 & 正态输出 $\in$ [0.90,\allowbreak{}0.99]；重尾输出低于它(\allowbreak{}差距在二项标准误 0.013 的 2 倍量级) & 正态 0.9533(\allowbreak{}命中 286/\allowbreak{}300)vs 重尾 0.9267(\allowbreak{}命中 278/\allowbreak{}300)，差 +0.0267(\allowbreak{}se$\approx$0.0122)，半宽 0.2781 vs 0.2856 \checkmark \\
C & N=1000 时同一个重尾输出的覆盖率(差距是否消失) & 两者都回到 0.95 附近(\allowbreak{}|覆盖率 - 0.95| <=\allowbreak{} 3 se) & 正态 0.9467，重尾 0.9567，|重尾-0.95|=\allowbreak{}0.0067(\allowbreak{}3se=\allowbreak{}0.0353) \checkmark \\
B & 精确参考下的 KS 序列 / CLT 参考下的 KS 序列 & 精确参考：KS(\allowbreak{}32000) < KS(\allowbreak{}2000)；CLT 参考：KS 不随 N 消失(\allowbreak{}地板 > 5e-3) & 精确 [0.01492 0.00722 0.00422](\allowbreak{}N=\allowbreak{}(\allowbreak{}2000,\allowbreak{} 8000,\allowbreak{} 32000))，CLT [0.0265  0.02037 0.01766] \checkmark \\
A & Sobol 网三种置乱下的每格点数极值 vs MC & Sobol：min=\allowbreak{}max=\allowbreak{}1(\allowbreak{}所有 m1+m2=\allowbreak{}m 的切法)；MC：min=\allowbreak{}0 且 max >=\allowbreak{} 2 & m=\allowbreak{}4/\allowbreak{}none: [1,\allowbreak{}1]；m=\allowbreak{}4/\allowbreak{}shift: [1,\allowbreak{}1]；m=\allowbreak{}4/\allowbreak{}linear: [1,\allowbreak{}1]；m=\allowbreak{}6/\allowbreak{}none: [1,\allowbreak{}1]；m=\allowbreak{}6/\allowbreak{}shift: [1,\allowbreak{}1]；m=\allowbreak{}6/\allowbreak{}linear: [1,\allowbreak{}1]；m=\allowbreak{}8/\allowbreak{}none: [1,\allowbreak{}1]；m=\allowbreak{}8/\allowbreak{}shift: [1,\allowbreak{}1]；m=\allowbreak{}8/\allowbreak{}linear: [1,\allowbreak{}1]；MC 256 点在 16x16 格上 [0,\allowbreak{}4](\allowbreak{}空格 86 个) \checkmark \\
A & HK 不等式被违反的次数 / 松紧度的分布 & 违反 0 次；松紧度 min < 1e-2 且 max > 0.5 & 违反 0/\allowbreak{}20 次；松紧度 min=\allowbreak{}0.0000，max=\allowbreak{}0.9685，中位数=\allowbreak{}0.0288(\allowbreak{}跨 300.0 个数量级) \checkmark \\
A & 变差的频率标度 V(k)/(V(1)k\textasciicircum{}2)(两种网格分辨率) & =\allowbreak{} 1；400 格上偏差 <=\allowbreak{} 2e-3(\allowbreak{}差分离散化)，1600 格上再小 4 倍以上 & V(\allowbreak{}2)/\allowbreak{}(\allowbreak{}V(\allowbreak{}1)*2\textasciicircum{}2)=\allowbreak{}0.999814466；V(\allowbreak{}4)/\allowbreak{}(\allowbreak{}V(\allowbreak{}1)*4\textasciicircum{}2)=\allowbreak{}0.999072560；V(\allowbreak{}8)/\allowbreak{}(\allowbreak{}V(\allowbreak{}1)*8\textasciicircum{}2)=\allowbreak{}1.000067037；400 格最大偏差 9.27e-04，1600 格最大偏差 2.43e-04 \checkmark \\
A & D*(37) 的 MC/Sobol 序列与 D*(64) 的地板 & 每个 N 上 D*(\allowbreak{}Sobol) < D*(\allowbreak{}MC)/\allowbreak{}4，且优势随 N 增大；D*(\allowbreak{}64) =\allowbreak{} 0(\allowbreak{}N=\allowbreak{}4096 时，锚网格分辨率地板) & 37 格：MC [0.06153 0.03517 0.01768]，Sobol [0.01152 0.00275 0.00087]；64 格 Sobol [0.00562 0.00073 0.     ] \checkmark \\
B & 确定性复制的标准差 / 真实误差 / 覆盖率 / 随机化后的误差棒 & 标准差 =\allowbreak{} 0、\allowbreak{}误差 > 1e-4、\allowbreak{}覆盖率 =\allowbreak{} 0、\allowbreak{}随机化后 stderr > 0 且 bar\_honesty*sqrt(\allowbreak{}R) $\approx$ 1 & 复制标准差 0.0e+00，误差 2.866e-04，确定性覆盖率 0.00，随机化 stderr\_mean 6.997e-07，bar\_honesty*sqrt(\allowbreak{}R)=\allowbreak{}0.986 \checkmark \\
B & 批量均值与真实标准误的比值(三种方法) & MC $\in$ [0.75,\allowbreak{}1.35]、\allowbreak{}Sobol $\in$ [0.5,\allowbreak{}1.5]、\allowbreak{}格点 $\in$ [0.5,\allowbreak{}1.5] & mc: 批量均值 2.773e-04 /\allowbreak{} 真实标准误 2.426e-04 =\allowbreak{} 1.143；sobol: 批量均值 2.390e-07 /\allowbreak{} 真实标准误 2.952e-07 =\allowbreak{} 0.810；lattice: 批量均值 2.941e-05 /\allowbreak{} 真实标准误 3.084e-05 =\allowbreak{} 0.954 \checkmark \\
C & 小窗口(帧所用)与大窗口上的 MC / Sobol 斜率 & 小窗口：MC $\in$ [-0.58,\allowbreak{}-0.42] 且 Sobol 至少快 0.30；大窗口：Sobol 更陡(\allowbreak{}斜率随窗口变) & 小窗口 [64,\allowbreak{}4096]：MC -0.4969，Sobol -1.2185(\allowbreak{}差 +0.7216)；大窗口 [1024,\allowbreak{}65536]：MC -0.5151，Sobol -1.9034 \checkmark \\
C & 间断边界在 2\textasciicircum{}-1(与网对齐)vs 在 0.3(不对齐) & 对齐：14 个 (\allowbreak{}N,\allowbreak{} 随机化) 组合的误差全为 0；不对齐：斜率 $\approx$ -0.85 且误差非零 & 对齐误差 [0. 0. 0. 0. 0. 0. 0. 0. 0. 0. 0. 0. 0. 0.](\allowbreak{}max=\allowbreak{}0e+00)；不对齐斜率 -0.898，误差 [0.01 0.01 0.   0.   0.   0.   0.  ] \checkmark \\
A & 生成矩阵累加 vs sobol\_\allowbreak{}points 的集合差 / 置乱后的 1 维分层 & 集合差 =\allowbreak{} 0；置乱后 256 点落在 256 个 1 维格子里，每格恰好 1 个 & 集合差 0.000e+00；置乱后 1 维 256 格(\allowbreak{}n=\allowbreak{}256)的计数极值 [1,\allowbreak{}1] \checkmark \\
A & Ishigami 解析值 vs 闭式公式 & 相对差 <=\allowbreak{} 1e-12；sum S\_i =\allowbreak{} 0.7563 < 1 < sum S\_Ti =\allowbreak{} 1.2437 & 最大相对差 2.776e-17；S=\allowbreak{}[0.3139 0.4424 0.    ]，S\_T=\allowbreak{}[0.5576 0.4424 0.2437]，sum S\_i=\allowbreak{}0.7563，sum S\_Ti=\allowbreak{}1.2437，V=\allowbreak{}13.8446 \checkmark \\
A & g-函数解析值 vs 闭式 + 错误公式的对照 & 正确公式相对差 <=\allowbreak{} 1e-12 且 sum S\_i < 1；错误公式 sum S\_i =\allowbreak{} 1.0716 > 1 & 最大相对差 0.000e+00；正确 sum S\_i=\allowbreak{}0.9290，sum S\_Ti=\allowbreak{}1.0716；错误公式 sum S\_i=\allowbreak{}1.0716 \checkmark \\
A & Gram 矩阵偏离单位阵 / 节点与权重对拍 / Hermite 三点权重 & Gram <=\allowbreak{} 1e-9；节点 <=\allowbreak{} 5e-15；权重比恒为 2(\allowbreak{}<=\allowbreak{} 1e-15)；Hermite 三点 =\allowbreak{} (\allowbreak{}1/\allowbreak{}6,\allowbreak{}2/\allowbreak{}3,\allowbreak{}1/\allowbreak{}6) & Gram 最大偏离 4.944e-11；GL 节点 1.887e-15；GL 权重差 6.731e-16；GH 三点权重 [0.166667 0.666667 0.166667] \checkmark \\
A & poly4 在 p=2,3,4,6,12 的投影误差 与 |x| 在 p=12 的误差 & poly4：p<4 时 > 1e-2、\allowbreak{}p>=\allowbreak{}4 时 <=\allowbreak{} 1e-12；|x| 在 p=\allowbreak{}12 时 > 1e-2 & poly4 [1.292e-01 9.158e-02 2.457e-15 4.020e-15 1.627e-14]；|x| p=\allowbreak{}12 $\to$ 0.0319 \checkmark \\
A & PCE 解析指数 vs pick-freeze 实测(双线性，N=4096) & PCE 精确给出 (\allowbreak{}0.5,\allowbreak{}0.5)(\allowbreak{}相对 1e-12)；pick-freeze 在 $\pm$0.05 内 & PCE S=\allowbreak{}[0.5 0.5]，S\_T=\allowbreak{}[0.5 0.5]；pick-freeze S=\allowbreak{}[0.5046 0.5064](\allowbreak{}半宽 [0.0315 0.0306])，V=\allowbreak{}2.0016 \checkmark \\
B & 独立输入：g-函数的 sum S\_\allowbreak{}i / sum S\_\allowbreak{}Ti 与解析值、区间 & sum S\_i - 半宽 <=\allowbreak{} 1(\allowbreak{}不显著超过 1)、\allowbreak{}sum S\_Ti >=\allowbreak{} 1(\allowbreak{}区间内)，且与解析值差 <=\allowbreak{} 0.08 & sum S\_i=\allowbreak{}0.9348$\pm$0.0820(\allowbreak{}解析 0.9290)，sum S\_Ti=\allowbreak{}1.1077$\pm$0.0358(\allowbreak{}解析 1.0716) \checkmark \\
B & 相关输入(rho=0.5)：条件方差口径 / 乘积测度口径 & 条件方差口径 > 1(\allowbreak{}实测 1.4942，超出部分 0.4942)；乘积测度口径 $\approx$ 1(\allowbreak{}实测 0.989) & 条件方差口径 1.4942(\allowbreak{}超出 +0.4942)，乘积测度口径 0.9890，两者差 +0.5052 \checkmark \\
B & sum S\_\allowbreak{}i(N) 的漂移幅度与解析值对照 & 漂移 >=\allowbreak{} 0.15；N=\allowbreak{}4096 处与解析值差 <=\allowbreak{} 0.08 & sum S\_i(\allowbreak{}N)=\allowbreak{}[1.1587 0.8416 1.1657 1.0034 0.8684 0.932  0.9488](\allowbreak{}N=\allowbreak{}[64,\allowbreak{} 128,\allowbreak{} 256,\allowbreak{} 512,\allowbreak{} 1024,\allowbreak{} 2048,\allowbreak{} 4096])，漂移 0.3242，解析 0.9290，末端差 0.0198 \checkmark \\
C & 两种被积函数的半对数衰减率 & exp：斜率 <=\allowbreak{} -0.9 且 R$^2$ >=\allowbreak{} 0.98；|x|：斜率 >=\allowbreak{} -0.2(\allowbreak{}远慢于几何) & exp: 半对数斜率 -1.0005(\allowbreak{}R$^2$=\allowbreak{}0.9919)；abs: 半对数斜率 -0.1080(\allowbreak{}R$^2$=\allowbreak{}0.8596)；exp 从 p=\allowbreak{}1 到 p=\allowbreak{}12 掉了 10.9 个数量级，|x| 只掉了 1.4 个数量级 \checkmark \\
C & |x| 的代数衰减指数(三种子样本) & 三者都落在 [0.9,\allowbreak{}1.5] 且 R$^2$ >=\allowbreak{} 0.97：与渐近的 p\textasciicircum{}-1(\allowbreak{}c\_k \textasciitilde{} k\textasciicircum{}(\allowbreak{}-3/\allowbreak{}2))一致，但 p<=\allowbreak{}16 上测不出干净的 -1(\allowbreak{}奇偶台阶) & 全部阶 alpha=\allowbreak{}1.2439(\allowbreak{}R$^2$=\allowbreak{}0.9811)；奇数阶 alpha=\allowbreak{}1.3061(\allowbreak{}R$^2$=\allowbreak{}0.9996)；偶数阶 alpha=\allowbreak{}1.0940(\allowbreak{}R$^2$=\allowbreak{}0.9972)；误差序列 [0.5099 0.1282 0.1282 0.0638 0.0646 0.0405 0.0405 0.0296 0.0297 0.0217
 0.022  0.0182 0.0188 0.0149 0.0149 0.0144] \checkmark \\
\end{longtable}

